\chapter{Vocabulary}

This chapter defines the manual's CharlotteOS-specific terms.

\section{Protection domain}

A \textbf{protection domain} is an address space and authority boundary. It owns:

\begin{itemize}
    \item A capability table.
    \item Memory mappings.
    \item One or more execution contexts (threads).
    \item Resource accounting.
    \item Failure and teardown state.
\end{itemize}

A \emph{process} is the concrete implementation of a protection domain, but
architectural discussion prefers the semantic term.

\section{Execution context}

An \textbf{execution context} is a kernel-scheduled thread. It may have:

\begin{itemize}
    \item LP affinity.
    \item Priority or scheduling class.
    \item A current protection domain.
    \item A wait state.
    \item An associated userspace executor.
\end{itemize}

\section{Logical processor (LP)}

A \textbf{logical processor} is CharlotteOS's representation of a schedulable
hardware execution context --- typically one CPU core (or one hyperthread).

An LP is \textbf{not} a shard (in the sitas sense). A typical mapping is:

\begin{verbatim}
one sitas shard
    <-> one LP-affine userspace thread
    <-> normally one logical CPU
\end{verbatim}

The architecture retains the distinction so that CPU hotplug, 
oversubscription, migration, asymmetric CPUs, and multiple sharded processes
remain possible.

\section{Sitas shard}

A \textbf{shard} is a userspace ownership and execution domain:

\begin{itemize}
    \item One executor runs its tasks.
    \item At most one task at a time mutates shard-owned state.
    \item Cross-shard access occurs through typed messages carrying owned values.
    \item Placement is explicit but is not the shard's identity.
\end{itemize}

\section{Capability}

A \textbf{capability} is an unforgeable reference to a kernel object, carrying
rights. It answers: ``May I perform this operation?''

Capabilities are scoped to a protection domain's capability table. A process
cannot fabricate a capability; it can only receive one through the bootstrap
contract or through IPC.

\section{Endpoint}

An \textbf{endpoint} is a bounded server-side IPC receive object. It has:

\begin{itemize}
    \item An owning protection domain.
    \item A queue capacity.
    \item A protocol/interface identity.
    \item Receive rights.
    \item Lifecycle and closure state.
\end{itemize}

\section{Connection capability}

A \textbf{connection capability} is a client-side authority to send or call an
endpoint. Possession of a service name is not authority. A name service resolves a
name and conditionally returns a connection capability.

\section{Message}

A \textbf{message} is a kernel-validated IPC envelope consisting of:

\begin{itemize}
    \item Interface or protocol identifier.
    \item Opcode.
    \item Flags.
    \item Inline scalar words.
    \item Zero or more memory-object attachments.
    \item Zero or more delegated capability attachments.
    \item Optional reply authority.
\end{itemize}

The kernel understands the envelope structure, not arbitrary Rust types.

\section{Reply token}

A \textbf{reply token} is one-shot authority to complete a blocking or deferred
call. It is:

\begin{itemize}
    \item Created by the kernel.
    \item Bound to the caller's pending call.
    \item Consumable exactly once.
    \item Invalidated by cancellation, caller death, or endpoint teardown.
    \item Optionally delegable when explicitly allowed.
\end{itemize}

\section{Resource capability}

A \textbf{resource capability} authorizes operations on a kernel-managed object
such as a memory object, endpoint, connection, interrupt, MMIO region, DMA domain,
device, completion queue, or timer.

\section{Operation and completion queue}

An \textbf{operation} is submitted work whose progress may depend on the kernel,
hardware, or a privileged service.

A \textbf{completion queue (CQ)} is the aggregated, per-shard or per-process
destination for operation results. It is waitable (with \texttt{CQ\_WAIT}) and has
non-lossy terminal completion semantics.

\section{Notification}

A \textbf{notification} means: ``State may have changed; inspect the corresponding
object or queue.'' It may be coalesced. It is not itself a result.

\section{Rust Waker}

A Rust \texttt{Waker} means: ``Poll this userspace future again.'' It is a
userspace scheduling hint. It is \textbf{not} a kernel completion record and
should not cross the ABI directly.

\section{Distinctions that matter}

Several concepts sound similar but must be kept separate. Confusing them leads to
architectural mistakes.

\begin{center}
\begin{tabular}{|l|l|l|}
\hline
\textbf{Concept} & \textbf{Is} & \textbf{Is not} \\
\hline
Capability & Unforgeable authority to act & An event, a file descriptor, a handle \\
Endpoint & Server-side bounded receive queue & A socket, a pipe, a channel \\
Completion & A result of an async operation & A Waker, a notification \\
Waker & A hint to repoll a future & A kernel completion record \\
Shard & A userspace ownership domain & An LP, a core, a thread \\
LP & A schedulable hardware context & A shard \\
Notification & ``Inspect state'' & A result, a data payload \\
IPC call & Cross-domain service invocation & A kernel operation submission \\
\hline
\end{tabular}
\end{center}

\section{Quick reference: acronyms}

\begin{description}
    \item[ASID] Address Space Identifier. 0 = kernel; $>$0 = user address space.
    \item[CQ] Completion Queue. Shared-memory ring buffer for async operation
        results.
    \item[EL0] Exception Level 0 on AArch64. CharlotteOS also retains
        \texttt{EL0} as an architecture-neutral source and test label for
        userspace execution, including ring 3 on x86-64.
    \item[EL1] Exception Level 1 (AArch64). Kernel execution.
    \item[HHDM] Higher Half Direct Mapping. Linear mapping of all physical memory
        into kernel virtual address space (provided by the Limine bootloader).
    \item[IPI] Inter-Processor Interrupt. Signal from one LP to another.
    \item[LP] Logical Processor. One schedulable hardware execution context.
    \item[MMU] Memory Management Unit. Hardware that enforces page-table
        permissions.
    \item[IOMMU / SMMU] I/O Memory Management Unit. Enforces DMA access
        permissions from devices, analogous to the MMU for CPU access.
\end{description}

\endinput
